\chapter{Discussion}
\label{chap:discussion}

\section{Interpreting the Three-Level Framework}
\label{sec:disc-framework-interpretation}
% What the three methods together reveal that no single method could show alone
% Cases where the levels align vs. diverge -- and what each divergence type means
% The framework as a diagnostic tool: which layer is the source of instability?
% [FIG] Stability failure taxonomy -- which divergence pattern indicates which
%   type of clinical risk

\section{Cross-Model Comparison}
\label{sec:disc-models}
% Mistral Large 3 vs. GPT-5.4 and open-weight peers across all three methods
% Sovereignty dimension: does the EU flagship match proprietary models?
% Does self-hostability come at a stability cost? -- the sovereignty tradeoff
% Size-class effects: small (24-32B) vs. mid (70B) vs. large MoE (~40B active)
%   vs. proprietary frontier. Does more parameters mean more stability?
% [GRAPH] Summary: all models x all three metrics x both temperatures --
%   the central comparative result

\section{Temperature and Clinical Deployment}
\label{sec:disc-temperature}
% Greedy decoding as an upper-bound baseline -- theoretical best case
% Stochastic drift patterns: which models degrade most under T=0.7?
% Clinical significance of measured drift patterns
% Practical deployment recommendations: when is T=0.7 acceptable?
% [GRAPH] Delta chart: T=0.7 minus T=0.0 per model -- the stochastic cost

\section{Vignette-Dependent Behaviour}
\label{sec:disc-vignettes}
% Which patient profiles challenge model consistency most: resistant, trauma
% Strategy selection patterns across clinical presentations -- are some
%   strategies inherently harder to maintain consistently?
% Implications for deployment: not all patient profiles are equally safe
% Implications for patient-specific adaptation
% [GRAPH] Per-vignette stability ranking across all three methods

\section{Evaluation Framework Reflections}
\label{sec:disc-framework-design}
% Strengths: hierarchical decomposition, reproducibility, clinical grounding
% What the taxonomy design choices reveal -- and what they constrain
%   Taxonomy sensitivity: how category design shapes measured consistency
% What declared intent can and cannot claim about model cognition
%   Plan-as-declared-intent: what this framing enables and what it cannot claim
% Generalisability: can this framework apply to CBT, DBT, other manualized protocols?
% [FIG] Generalisation sketch -- framework mapped onto a second structured
%   protocol (e.g. CBT thought records) to show reusability

\section{Toward In Silico Clinical Trials}
\label{sec:disc-silico}
% This framework as a reusable, reproducible validation template
% Regulatory relevance: what stability testing can and cannot certify
% The remaining gap: from benchmark stability to demonstrated clinical safety
% [FIG] Positioning diagram -- where this framework sits in a full clinical
%   validation pipeline; what it certifies and what remains open
